\documentclass[10pt]{article}

% ---------- Document Identity (update these only) ----------
\newcommand{\DocStage}{}
\newcommand{\DocTitle}{Your Road to Tier One SOF}
\newcommand{\ProgramName}{182-Week Civilian-to-Selection Readiness Pipeline}
\newcommand{\ProgramSubtitle}{}

% ---------- Page ----------
\usepackage[legalpaper,left=0.5in,right=0.5in,top=0.6in,bottom=0.45in]{geometry}

% ---------- Typography ----------
\usepackage[T1]{fontenc}
\usepackage[utf8]{inputenc}
\usepackage{libertinus}
\usepackage{microtype}
\usepackage{parskip}
\usepackage{ragged2e}
\usepackage{textcomp}
\AtBeginDocument{\color{black}}
\AtBeginDocument{\pagecolor{white}}
\AtBeginDocument{\justifying}

% ---------- Landscape pages ----------
\usepackage{pdflscape}

% ---------- Color ----------
\usepackage[table]{xcolor}
\definecolor{StageBlue}{HTML}{1F4E79}
\definecolor{StageBlueDark}{HTML}{163A5A}
\definecolor{LightGray}{HTML}{F2F4F7}
\definecolor{MidGray}{HTML}{D9DEE5}
\definecolor{RuleGray}{HTML}{C7CED8}
\definecolor{HeaderInk}{HTML}{000000}
\definecolor{Ink}{HTML}{000000}

% ---------- Utility ----------
\usepackage{enumitem}
\sloppy
\setlength{\emergencystretch}{2em}

% ---------- Boxes / UI ----------
\usepackage[most]{tcolorbox}

\tcbset{
  charterTitle/.style={
    enhanced,
    colback=StageBlue!4,
    colframe=StageBlue,
    boxrule=1.25pt,
    arc=3pt,
    left=16pt,right=16pt,top=14pt,bottom=14pt,
    borderline north={3.2pt}{0pt}{StageBlue},
    borderline west={4pt}{0pt}{StageBlue},
    borderline south={1.25pt}{0pt}{StageBlue},
    drop shadow={black!25!white},
  },
  charterRule/.style={
    enhanced,
    colback=LightGray,
    colframe=RuleGray,
    boxrule=0.85pt,
    arc=3pt,
    left=12pt,right=12pt,top=10pt,bottom=10pt,
    borderline west={3pt}{0pt}{StageBlue},
  },
  charterBadge/.style={
    enhanced,
    colback=StageBlue,
    coltext=white,
    colframe=StageBlueDark,
    boxrule=0.6pt,
    arc=2pt,
    left=6pt,right=6pt,top=3pt,bottom=3pt,
  }
}
\newtcbox{\Badge}{nobeforeafter,charterBadge}

% ---------- Lists ----------
\setlist[itemize]{leftmargin=1.5em, itemsep=0.25em, topsep=0.25em}
\setlist[enumerate]{leftmargin=1.8em, itemsep=0.25em, topsep=0.25em}

% ---------- TOC ----------
\usepackage{tocloft}
\setcounter{tocdepth}{3}
\setlength{\cftbeforesecskip}{3pt}
\setlength{\cftbeforesubsecskip}{2pt}
\setlength{\cftbeforesubsubsecskip}{0pt}
\renewcommand{\cftsecfont}{\bfseries}
\renewcommand{\cftsecpagefont}{\bfseries}
\renewcommand{\cftsubsecfont}{\normalfont}
\renewcommand{\cftsubsecpagefont}{\normalfont}
\renewcommand{\cftsubsubsecfont}{\normalfont}
\renewcommand{\cftsubsubsecpagefont}{\normalfont}
\setlength{\cftsecindent}{0pt}
\setlength{\cftsubsecindent}{1.25em}
\setlength{\cftsubsubsecindent}{2.8em}
\setlength{\cftsecnumwidth}{2.1em}
\setlength{\cftsubsecnumwidth}{2.8em}
\setlength{\cftsubsubsecnumwidth}{3.6em}

% Remove the built-in TOC heading so only the custom heading is shown
\makeatletter
\renewcommand{\tableofcontents}{\@starttoc{toc}}
\makeatother

% ---------- Header/Footer: page number only ----------
\usepackage{fancyhdr}
\pagestyle{fancy}
\fancyhf{}
\renewcommand{\headrulewidth}{0pt}
\renewcommand{\footrulewidth}{0pt}
\fancyfoot[C]{\thepage}

% ---------- Section formatting ----------
\usepackage{titlesec}

\titleformat{\section}
  {\Large\bfseries\color{Ink}}
  {\thesection}{0.6em}{}
  [\vspace{0.25em}\textcolor{StageBlue}{\rule{\linewidth}{0.8pt}}]

\titleformat{\subsection}
  {\large\bfseries\color{Ink}}
  {\thesubsection}{0.6em}{}

\titleformat{\subsubsection}
  {\normalsize\bfseries\color{Ink}}
  {\thesubsubsection}{0.6em}{}

\titlespacing*{\section}{0pt}{0.95em}{0.35em}
\titlespacing*{\subsection}{0pt}{0.65em}{0.25em}
\titlespacing*{\subsubsection}{0pt}{0.55em}{0.2em}

% ---------- Table engine ----------
\usepackage{tabularray}
\UseTblrLibrary{booktabs}

% ---------- tabularray: GLOBAL table treatment ----------
\SetTblrInner{
  cells = {fg=black, bg=white, valign=t},
}

% ---------- Header helpers ----------
\newcommand{\Hdr}[1]{\shortstack[c]{\strut #1\strut}}

% ---------- Lines ----------
\newcommand{\TitleLine}{\textcolor{StageBlue}{\rule{0.98\linewidth}{0.95pt}}}
\newcommand{\SubTitleLine}{\textcolor{RuleGray}{\rule{0.98\linewidth}{0.65pt}}}

% ---------- Continued on next page ----------
\newcommand{\ContinuedOnNextPageInline}{%
  \par
  \vfill
  \noindent\textcolor{RuleGray}{\rule{\linewidth}{1pt}}\vspace{-2pt}
  \noindent\hfill\textit{Continued on next page}\par\vspace{-10pt}
  \noindent\textcolor{RuleGray}{\rule{\linewidth}{1pt}}
  \clearpage
}

% ---------- PDF hyperlinks ----------
\usepackage[hidelinks]{hyperref}
\hypersetup{
  colorlinks=false,
  pdfborder={0 0 0},
  linktoc=all,
  pdftitle={\DocTitle},
  pdfauthor={\ProgramName}
}

% =========================================================
\begin{document}

\section*{Stage 3.D — Integrated Phase Calendars Compact Form}
\phantomsection
\addcontentsline{toc}{section}{Stage 3.D — Integrated Phase Calendars Compact Form}

\subsection*{3.D.1 Calendar Encoding and Use Rules (Binding)}
\phantomsection
\addcontentsline{toc}{subsection}{3.D.1 Calendar Encoding and Use Rules (Binding)}

\begin{enumerate}
\item This deliverable provides compact, phase-level calendars. It defines protected window placement by week-in-phase and identifies the intended major test constructs inside those protected windows. It does not prescribe day-level schedules, workouts, sets/reps, pacing, or microcycle design.

\item “Deload+Test” weeks are protected windows by default doctrine. They are the authorized weeks for gate-grade test execution when validity, safety, and comparability prerequisites are satisfied. “Build” and “Intensify” weeks are execution weeks where major gate-grade tests are not scheduled by default; low-cost monitoring continues.

\item Stage 2 is authoritative for test definitions, protocol IDs, metric IDs, validity fields, Evidence Ref rules, and environment/tool identity posture. Stage 3 is authoritative for placement and reslot posture. Any mismatch must be reconciled in favor of Stage 2 identifiers and Stage 1.F admissibility posture.

\item Major test event caps and spacing doctrine apply inside each Deload+Test week (protected window):
\begin{itemize}
\item Ideal cap: 2 major test events per Deload+Test week.
\item Hard cap: 3 major test events per Deload+Test week.
\item Any two major test events \textbf{MUST} NOT start within the same 72-hour window (start-to-start).
\item Maximal \textbf{RUN} and maximal \textbf{RUCK} \textbf{MUST} NOT start within the same 96-hour window (start-to-start).
\item If caps/spacing cannot be maintained inside the Deload+Test week, the correct posture is partial completion and reslot of remaining major tests to the next protected window. Compression is prohibited.
\end{itemize}

\item Major test events (count toward cap/spacing when executed as formal tests):
\begin{itemize}
\item \textbf{CAL} battery: \textbf{C4} (push-ups), \textbf{C5} (sit-ups/curl-ups), \textbf{C6} (pull-ups), \textbf{C7} (plank) executed as a standardized battery may be treated as one major test event for cap/spacing purposes.
\item \textbf{RUN} time trials: \textbf{C8} (1.5 mi), \textbf{C9} (2 mi), \textbf{C10} (3 mi), \textbf{C11} (5 mi).
\item \textbf{RUCK} time trials: \textbf{C12} (distance/load variants by \textbf{MET-2B-018}–021).
\item \textbf{SWIM} time trials: \textbf{C13} (500 yd / 500 m, unit locked).
\item \textbf{STR} proxy testing: \textbf{C17} (submax proxy set; \textbf{MET-2B-027}–031).
\item \textbf{AER} time trials: \textbf{C18} (rower/bike/elliptical time trials; \textbf{MET-2B-032}–034).
\end{itemize}

\item Low-cost monitoring does not count toward major test caps and may occur in any week:
\begin{itemize}
\item \textbf{BODYCOMP} trend capture (\textbf{C1}–\textbf{C3}) at the required cadence.
\item \textbf{RECOVERY} marker capture (\textbf{C15}) daily/per session.
\item \textbf{DURABILITY} screen flags and pain/foot tracking (\textbf{C16}) daily/per session.
\end{itemize}

\item Environment and tool identity are mandatory for gate-grade tests:
\begin{itemize}
\item \textbf{RUN}/\textbf{RUCK} route or machine identity must use Environment Unit IDs (\textbf{ROUTE}-2E-\#\#\# / \textbf{TM}-2E-\#\#\#) and must meet Stage 2 logging requirements (grade settings for treadmill; route verification posture).
\item \textbf{SWIM} requires \textbf{POOL}-2E-\#\#\# and verified unit/length; yards and meters are separate evidence streams.
\item Major tool changes (timer method, footwear, ruck configuration, pool unit/length) are comparability breaks and bias toward reslot until stabilized and replicated under the new stable configuration.
\end{itemize}

\item Underwater (\textbf{C14}) is conditional and governance-gated:
\begin{itemize}
\item Scheduled only in Phase 10–Phase 13 and only when certified lifeguard plus dedicated safety spotter and facility permission prerequisites are satisfied and logged.
\item If prerequisites are not met, underwater is not scheduled and is not attempted; no proxy substitutions are used for underwater readiness claims.
\end{itemize}

\item “Approx Program Week Range” is a convenience index computed from the frozen phase-duration ledger for the 182-week horizon. Phase-relative weeks (Week-in-Phase) are the controlling schedule coordinates.
\end{enumerate}

\newpage

\begin{landscape}

\subsection*{Phase 00 — Bodily Reactivation and Baseline Creation}
\phantomsection
\addcontentsline{toc}{subsubsection}{Phase 00 — Bodily Reactivation and Baseline Creation}

\begingroup
\scriptsize
\setlength{\tabcolsep}{2pt}
\renewcommand{\arraystretch}{1.12}

\begin{longtblr}{
  width=\linewidth,
  rowhead=1,
  colspec = {
    Q[l,wd=.7in]
    Q[l,wd=.80in]
    Q[l,wd=2.05in]
    X[l,1.25]
    X[l,3.50]
  },
  hlines,
  vlines,
  cells = {hsep=6pt, vsep=6pt, halign=l, valign=t},
  row{1} = {bg=StageBlue!15, fg=black, font=\bfseries, halign=l, valign=m},
  row{odd} = {bg=white},
  row{even} = {bg=LightGray},
}
\Hdr{Week-in-Phase} &
\Hdr{Week Range} &
\Hdr{Stress Type} &
\Hdr{Key Tests This Week (Major Tests Only)} &
\Hdr{Notes / Adjustments (Validity, Caps/Spacing, Environment, Reslot)} \\

Week 1 (baseline) &
Program W1 &
Deload+Test &
\textbf{BODYCOMP}: \textbf{C1} (\textbf{MET-2B-001}); \textbf{C2} (\textbf{MET-2B-002}); optional minimal \textbf{C3} (\textbf{MET-2B-007} only)
\textbf{CAL} minimal: \textbf{C7} (\textbf{MET-2B-013}) &
Baseline window is infrastructure-first: stable tool identity, stable Site \textbf{ID}, and clean Stage 2.E logging posture.
\textbf{CAL} is limited to plank as low-risk format/timing check.
Do not schedule major locomotion (\textbf{RUN}/\textbf{RUCK}) unless mechanics and safety posture are stable and environment is safe.
Major event count typically ≤1.
Any missing required fields make evidence inadmissible for gate use and drive reslot. \\

Weeks 2–3 &
Program W2–W3 &
Build &
None (major tests not scheduled) &
Training-focused.
Low-cost monitoring continues (\textbf{C15} \textbf{RECOVERY} daily; \textbf{C16} \textbf{DURABILITY} flags per session; \textbf{BODYCOMP} weekly).
No gate-grade \textbf{RUN}/\textbf{RUCK} tests outside protected windows.
Preserve environment/tool stability approaching Week 4. \\

Week 4 &
Program W4 &
Deload+Test &
Optional walk-based aerobic monitor only (Phase 00-local construct until reconciled) &
Phase 00 \textbf{RUN} is optional and conservative; do not force impact.
If outdoor footing/visibility is unsafe, treadmill is allowed only with \textbf{TM}-2E-\#\#\# and full settings logging.
If neither is feasible safely, reslot \textbf{RUN} to Week 8.
Major event count typically ≤1.
Do not introduce \textbf{CAL} battery here unless explicitly authorized later; plank-only remains default. \\

Weeks 5–7 &
Program W5–W7 &
Build &
None (major tests not scheduled) &
Training-focused.
Grand Forks winter hazards bias toward indoor substitutions for training exposures; do not attempt gate-grade tests outside windows.
Maintain stable footwear interface; footwear changes are logged and treated as comparability breaks near test windows. \\

Week 8 &
Program W8 &
Deload+Test &
Walk-based aerobic monitor check-in only (Phase 00-local construct until reconciled)
\textbf{BODYCOMP} minimal set continues as trend support &
\textbf{RUN} check-in remains capacity-focused, not performance-chasing.
Enforce Stage 2 validity fields and Evidence Ref.
If fatigue or durability flags are elevated (\textbf{C16} gait/foot flags; \textbf{C15} fatigue trend), treat the window as recovery-first and reslot \textbf{RUN} to Week 12 rather than forcing compromised evidence.
Major event cap ≤1. \\

Weeks 9–11 &
Program W9–W11 &
Intensify &
None (major tests not scheduled) &
Training-focused.
Preserve route/machine stability and timing method stability approaching Week 12.
Do not introduce new environment unit IDs immediately before the next protected window unless forced (if forced, log as comparability break and plan conservative interpretation). \\

Week 12 &
Program W12 &
Deload+Test &
\textbf{RUN} check-in: \textbf{C8} (\textbf{MET-2B-014}) (repeat under comparable conditions to support early replication posture when feasible) &
Replication enablement is optional in Phase 00; if only one \textbf{VALID} run exists in-window, do not force additional attempts in the same week.
Reslot to Week 16 if replication is needed and conditions can be stabilized.
Any mechanics-altering pain triggers \textbf{STOP} \textbf{NOW} and reslot.
Major event cap ≤1. \\

Weeks 13–15 &
Program W13–W15 &
Intensify &
None (major tests not scheduled) &
Training-focused.
Validity protection posture increases approaching end-of-phase.
Avoid novelty in footwear, timing method, and route selection in the last 2 weeks before Week 16. \\

Week 16 &
Program W16 &
Deload+Test &
End-of-phase conservative verification: walk-based aerobic monitor only (Phase 00-local construct until reconciled)
\textbf{CAL} minimal \textbf{C7} (\textbf{MET-2B-013}) if needed for replication &
End-of-phase window remains conservative.
Do not “upgrade” this window into a multi-domain battery.
If conditions are compromised (ice, wind, illness, missing logging), reslot rather than forcing evidence.
Treat this as an exit touchpoint for infrastructure readiness to support Phase 01’s more formal \textbf{CAL} battery. \\

\end{longtblr}

\endgroup

\newpage

\subsection*{Phase 01 — Base Conditioning and Hypertrophy}
\phantomsection
\addcontentsline{toc}{subsubsection}{Phase 01 — Base Conditioning and Hypertrophy}

\begingroup
\scriptsize
\setlength{\tabcolsep}{2pt}
\renewcommand{\arraystretch}{1.12}

\begin{longtblr}{
  width=\linewidth,
  rowhead=1,
  colspec = {
    Q[l,wd=.7in]
    Q[l,wd=.80in]
    Q[l,wd=2.05in]
    X[l,1.25]
    X[l,3.50]
  },
  hlines,
  vlines,
  cells = {hsep=6pt, vsep=6pt, halign=l, valign=t},
  row{1} = {bg=StageBlue!15, fg=black, font=\bfseries, halign=l, valign=m},
  row{odd} = {bg=white},
  row{even} = {bg=LightGray},
}
\Hdr{Week-in-Phase} &
\Hdr{Week Range} &
\Hdr{Stress Type} &
\Hdr{Key Tests This Week (Major Tests Only)} &
\Hdr{Notes / Adjustments (Validity, Caps/Spacing, Environment, Reslot)} \\

Week 1 (baseline) &
Program W17 &
Deload+Test &
\textbf{BODYCOMP}: \textbf{C1} (\textbf{MET-2B-001}); \textbf{C2} (\textbf{MET-2B-002}); \textbf{C3} minimal if used
\textbf{CAL} baseline battery intent: \textbf{C4}–\textbf{C7} (\textbf{MET-2B-010}–013) &
Baseline \textbf{CAL} battery intent exists to establish strict standards and artifact workflow.
\textbf{CAL} gate-validity requires compliant artifact evidence; missing/non-compliant artifact makes \textbf{CAL} tests \textbf{INVALID} for gate use and triggers reslot to Week 4.
Keep major event count at 1 (\textbf{CAL} battery) unless explicitly adding \textbf{RUN} is necessary and spacing can be preserved (usually not in baseline). \\

Weeks 2–3 &
Program W18–W19 &
Build &
None (major tests not scheduled) &
Training-focused.
Low-cost monitoring continues.
Protect tool/environment stability approaching Week 4: route identity, treadmill identity, footwear, timer method, and \textbf{CAL} artifact setup. \\

Week 4 &
Program W20 &
Deload+Test &
\textbf{CAL} battery: \textbf{C4}–\textbf{C7}
\textbf{RUN} \textbf{TT}: \textbf{C8} (\textbf{MET-2B-014}) &
Two major events.
Enforce 72-hour spacing between \textbf{CAL} battery and \textbf{RUN} \textbf{TT} within the week; if spacing cannot be preserved due to schedule constraints, reslot \textbf{RUN} to Week 8.
\textbf{RUN} is gate-relevant only if fully logged per Stage 2 (\textbf{ROUTE}-2E-\#\#\# or \textbf{TM}-2E-\#\#\#; treadmill grade logged). \\

Weeks 5–7 &
Program W21–W23 &
Intensify &
None (major tests not scheduled) &
Training-focused.
No gate-grade time trials outside protected windows.
Avoid novelty in footwear and routes approaching Week 8. \\

Week 8 &
Program W24 &
Deload+Test &
\textbf{CAL} battery: \textbf{C4}–\textbf{C7}
\textbf{RUN} \textbf{TT}: \textbf{C9} (\textbf{MET-2B-015}) &
Two major events.
Enforce 72-hour spacing.
If outdoor conditions compromise validity, treadmill substitution is permitted only with full settings logged.
Do not treat treadmill and outdoor results as interchangeable unless environment/tool stability is maintained across replication events. \\

Weeks 9–11 &
Program W25–W27 &
Intensify &
None (major tests not scheduled) &
Training-focused.
Preserve no-novelty posture approaching Week 12. \\

Week 12 &
Program W28 &
Deload+Test &
\textbf{CAL} battery: \textbf{C4}–\textbf{C7} (replication support window) &
Single major event by default to support \textbf{CAL} replication without overloading.
If \textbf{CAL} evidence is already replicated and stable, this window may be treated as recovery-first and used for conservative monitoring only; do not insert additional major tests unless replication gaps require it and spacing remains protected. \\

Weeks 13–15 &
Program W29–W31 &
Intensify &
None (major tests not scheduled) &
Training-focused. \\

Week 16 &
Program W32 &
Deload+Test &
\textbf{CAL} battery: \textbf{C4}–\textbf{C7}
\textbf{RUN} \textbf{TT}: \textbf{C9} (\textbf{MET-2B-015}) &
Two major events.
This window supports \textbf{RUN} replication in the 2-mile construct under stable conditions.
If replication status is still 1/3 for \textbf{RUN}, do not force immediate reattempts in the same week; plan next attempt in Week 20 or Week 22 per phase priorities and spacing rules. \\

Weeks 17–19 &
Program W33–W35 &
Intensify &
None (major tests not scheduled) &
Training-focused.
Avoid tool drift (new treadmill, new route, new timer method) immediately before Week 20. \\

Week 20 &
Program W36 &
Deload+Test &
\textbf{STR} proxy session: \textbf{C17}, only if finalized and deployable, (\textbf{MET-2B-027}-031)
\textbf{CAL} battery: \textbf{C4}–\textbf{C7} &
Two major events.
\textbf{STR} remains submax proxy only; treat as major event and enforce 72-hour spacing from \textbf{CAL} battery if both occur in the same week.
If spacing cannot be preserved, reslot \textbf{STR} to Week 22 or to the next phase baseline window.
Do not add \textbf{RUN} in Week 20 by default to avoid hard cap pressure. \\

Week 21 &
Program W37 &
Intensify &
None (major tests not scheduled) &
Training-focused.
This is a short runway into the end gate (Week 22).
No make-up stacking is permitted; preserve freshness and validity posture. \\

Week 22 (end gate) &
Program W38 &
Deload+Test &
\textbf{RUN} \textbf{TT}: \textbf{C9} (\textbf{MET-2B-015})
\textbf{CAL} battery: \textbf{C4}–\textbf{C7} &
End-of-phase exit gate window (22-week default).
Two major events.
Enforce 72-hour spacing and preserve stable environment/tool identity to support replication (2 of last 3 \textbf{VALID}).
If any required test is \textbf{INVALID}, default posture is reslot to the next available protected window under the chosen Phase 01 variant, or to Phase 02 Week 1 baseline. \\

\end{longtblr}

\endgroup

\paragraph*{Phase 01 variant note (24-week variant; 26-week extension):}

\begin{itemize}
\item 24-week variant: adds Weeks 23–24. Week 24 is a Deload+Test end gate window using the same major test set as Week 22. Program week ranges for all subsequent phases shift by +2 after Phase 01 under this variant.
\item 26-week extension: adds Weeks 25–26 as Transition and optional retest buffer only; no new major test types are introduced. Retests replace \textbf{INVALID} evidence or missing replication; they are not used to stack additional gate demands. Program week ranges for all subsequent phases shift by +4 after Phase 01 under this extension.
\end{itemize}

\newpage

\subsection*{Phase 02 — Maximal Strength Development (12 weeks; 3:1)}
\phantomsection
\addcontentsline{toc}{subsubsection}{Phase 02 — Maximal Strength Development}

\begingroup
\scriptsize
\setlength{\tabcolsep}{2pt}
\renewcommand{\arraystretch}{1.12}

\begin{longtblr}{
  width=\linewidth,
  rowhead=1,
  colspec = {
    Q[l,wd=.7in]
    Q[l,wd=.80in]
    Q[l,wd=2.05in]
    X[l,1.25]
    X[l,3.50]
  },
  hlines,
  vlines,
  cells = {hsep=6pt, vsep=6pt, halign=l, valign=t},
  row{1} = {bg=StageBlue!15, fg=black, font=\bfseries, halign=l, valign=m},
  row{odd} = {bg=white},
  row{even} = {bg=LightGray},
}
\Hdr{Week-in-Phase} &
\Hdr{Week Range} &
\Hdr{Stress Type} &
\Hdr{Key Tests This Week (Major Tests Only)} &
\Hdr{Notes / Adjustments (Validity, Caps/Spacing, Environment, Reslot)} \\

Week 1 (baseline) &
Program W39 &
Deload+Test &
\textbf{CAL} battery: \textbf{C4}–\textbf{C7};
\textbf{STR} proxy baseline: \textbf{C17}, only if finalized and deployable, (conservative) &
Baseline window introduces \textbf{STR} proxy under conservative submax posture.
If cap/spacing pressure exists, prioritize \textbf{CAL} battery and treat \textbf{STR} proxy as optional baseline marker.
Do not add \textbf{RUN} in baseline by default. \\

Weeks 2–3 &
Program W40–W41 &
Build &
None (major tests not scheduled) &
Training-focused; monitor \textbf{RECOVERY}/\textbf{DURABILITY} trends for readiness and to protect the Week 4 window. \\

Week 4 &
Program W42 &
Deload+Test &
\textbf{STR} proxy session: \textbf{C17}
\textbf{CAL} battery: \textbf{C4}–\textbf{C7} &
Two major events.
Enforce spacing within the week.
If mechanics-altering pain appears, terminate and reslot; do not force completion. \\

Weeks 5–7 &
Program W43–W45 &
Intensify &
None (major tests not scheduled) &
Training-focused; preserve environment/tool stability ahead of Week 8 \textbf{RUN} window. \\

Week 8 &
Program W46 &
Deload+Test &
\textbf{RUN} \textbf{TT}: \textbf{C9} (\textbf{MET-2B-015})
\textbf{CAL} battery: \textbf{C4}–\textbf{C7} &
Two major events.
Enforce 72-hour spacing.
If \textbf{RUN} is compromised by environment (ice/visibility), treadmill is allowed only under full \textbf{TM}-2E-\#\#\# logging.
Do not schedule \textbf{STR} proxy in Week 8 by default (would push toward hard cap and create spacing risk). \\

Weeks 9–11 &
Program W47–W49 &
Intensify &
None (major tests not scheduled) &
Training-focused; protect validity posture approaching Week 12 exit gate. \\

Week 12 &
Program W50 &
Deload+Test &
Exit gate intent: \textbf{RUN} \textbf{TT} \textbf{C9}
\textbf{CAL} battery
\textbf{STR} proxy session \textbf{C17} (staged; cap-limited) &
End-of-phase gate week may require up to 3 major events.
Execute only if hard cap and spacing can be satisfied inside the week.
Default prioritization if spacing cannot be protected: \textbf{RUN} + \textbf{CAL} as decisive; reslot \textbf{STR} to Phase 03 Week 1 baseline window.
Do not compress by placing two major events inside 72 hours.
Any \textbf{INVALID} test forces reslot posture. \\

\end{longtblr}

\endgroup

\newpage

\subsection*{Phase 03 — Converting Strength to Power and Endurance for SOF Selection}
\phantomsection
\addcontentsline{toc}{subsubsection}{Phase 03 — Converting Strength to Power and Endurance for SOF Selection}

\begingroup
\scriptsize
\setlength{\tabcolsep}{2pt}
\renewcommand{\arraystretch}{1.12}

\begin{longtblr}{
  width=\linewidth,
  rowhead=1,
  colspec = {
    Q[l,wd=.7in]
    Q[l,wd=.80in]
    Q[l,wd=2.05in]
    X[l,1.25]
    X[l,3.50]
  },
  hlines,
  vlines,
  cells = {hsep=6pt, vsep=6pt, halign=l, valign=t},
  row{1} = {bg=StageBlue!15, fg=black, font=\bfseries, halign=l, valign=m},
  row{odd} = {bg=white},
  row{even} = {bg=LightGray},
}
\Hdr{Week-in-Phase} &
\Hdr{Week Range} &
\Hdr{Stress Type} &
\Hdr{Key Tests This Week (Major Tests Only)} &
\Hdr{Notes / Adjustments (Validity, Caps/Spacing, Environment, Reslot)} \\

Week 1 (baseline) &
Program W51 &
Deload+Test &
\textbf{CAL} battery: \textbf{C4}–\textbf{C7}
Optional \textbf{RUN} check-in: \textbf{C8} (\textbf{MET-2B-014}) &
Baseline \textbf{RUN} is check-in only.
Do not treat as decisive gate by default.
Maintain low density: 1–2 major events maximum; if adding \textbf{RUN} would create spacing problems with \textbf{CAL}, reslot \textbf{RUN} to Week 4. \\

Weeks 2–3 &
Program W52–W53 &
Build &
None (major tests not scheduled) &
Training-focused. \\

Week 4 &
Program W54 &
Deload+Test &
\textbf{RUN} \textbf{TT}: \textbf{C8} (\textbf{MET-2B-014})
\textbf{CAL} battery: \textbf{C4}–\textbf{C7} &
Two major events.
Enforce 72-hour spacing.
If \textbf{RUN} mechanics are unstable or environment unsafe, reslot rather than forcing a compromised attempt. \\

Weeks 5–7 &
Program W55–W57 &
Intensify &
None (major tests not scheduled) &
Training-focused; preserve route/tool stability. \\

Week 8 &
Program W58 &
Deload+Test &
\textbf{RUN} \textbf{TT}: \textbf{C9} (\textbf{MET-2B-015})
\textbf{CAL} battery
Optional \textbf{STR} proxy \textbf{C17}, only if finalized and deployable, (cap-limited) &
Default is \textbf{RUN} + \textbf{CAL} (2 major events).
\textbf{STR} proxy is supporting and should be reslotted first if cap/spacing is threatened.
If three events are executed, enforce spacing within the week; do not cluster \textbf{RUN} and \textbf{STR} in a way that contaminates the \textbf{CAL} battery or produces durability breakdown signals. \\

Weeks 9–11 &
Program W59–W61 &
Intensify &
None (major tests not scheduled) &
Training-focused; protect validity posture approaching Week 12. \\

Week 12 &
Program W62 &
Deload+Test &
\textbf{RUN} \textbf{TT}: \textbf{C10} (\textbf{MET-2B-016})
\textbf{CAL} battery: \textbf{C4}–\textbf{C7} &
End-of-phase exit gate includes introduction of 3-mile \textbf{RUN} time trial.
Keep the week at 2 major events by default.
Any attempt to add \textbf{RUCK} in the same week must respect cap/spacing and should be avoided unless explicitly required by phase contract. \\

\end{longtblr}

\endgroup

\newpage

\subsection*{Phase 04 — Strength-Endurance and Work Capacity Development}
\phantomsection
\addcontentsline{toc}{subsubsection}{Phase 04 — Strength-Endurance and Work Capacity Development}

\begingroup
\scriptsize
\setlength{\tabcolsep}{2pt}
\renewcommand{\arraystretch}{1.12}

\begin{longtblr}{
  width=\linewidth,
  rowhead=1,
  colspec = {
    Q[l,wd=.7in]
    Q[l,wd=.80in]
    Q[l,wd=2.05in]
    X[l,1.25]
    X[l,3.50]
  },
  hlines,
  vlines,
  cells = {hsep=6pt, vsep=6pt, halign=l, valign=t},
  row{1} = {bg=StageBlue!15, fg=black, font=\bfseries, halign=l, valign=m},
  row{odd} = {bg=white},
  row{even} = {bg=LightGray},
}
\Hdr{Week-in-Phase} &
\Hdr{Week Range} &
\Hdr{Stress Type} &
\Hdr{Key Tests This Week (Major Tests Only)} &
\Hdr{Notes / Adjustments (Validity, Caps/Spacing, Environment, Reslot)} \\

Week 1 (baseline) &
Program W63 &
Deload+Test &
\textbf{CAL} battery: \textbf{C4}–\textbf{C7} &
Baseline window for Phase 04 remains conservative.
Keep major event count at 1 unless explicit replication gaps exist. \\

Weeks 2–3 &
Program W64–W65 &
Build &
None (major tests not scheduled) &
Training-focused; monitor \textbf{DURABILITY} foot/gait markers to determine readiness for ruck testing later in the phase. \\

Week 4 &
Program W66 &
Deload+Test &
\textbf{RUCK} \textbf{TT}: \textbf{C12} (\textbf{MET-2B-019}; 6 mi / 35 lb dry)
\textbf{CAL} battery (optional if cap allows) &
\textbf{RUCK} is high-cost; if \textbf{CAL} battery is included, enforce hard cap and spacing.
Default posture may be \textbf{RUCK} stand-alone (1 major event) if foot integrity is marginal or environment volatility is high.
Load must be verified (dry load); water logged separately.
Any foot hotspot trend forces reslot. \\

Weeks 5–7 &
Program W67–W69 &
Intensify &
None (major tests not scheduled) &
Training-focused; avoid novelty in ruck configuration and footwear approaching next \textbf{RUN} window. \\

Week 8 &
Program W70 &
Deload+Test &
\textbf{RUN} \textbf{TT}: \textbf{C10} (\textbf{MET-2B-016})
\textbf{CAL} battery &
Two major events.
Keep \textbf{RUCK} out of this window by default to protect spacing and durability.
\textbf{RUN} environment must be stable and logged; treadmill allowed only under full settings logging. \\

Weeks 9–11 &
Program W71–W73 &
Intensify &
None (major tests not scheduled) &
Training-focused; protect stability ahead of end-of-phase gate. \\

Week 12 &
Program W74 &
Deload+Test &
Exit gate intent:
\textbf{RUN} \textbf{TT} \textbf{C9} (\textbf{MET-2B-015})
\textbf{RUCK} \textbf{TT} \textbf{C12} (\textbf{MET-2B-019})
\textbf{CAL} battery (staged; cap-limited) &
End-of-phase may require up to 3 major events only if spacing is feasible.
Default prioritization if spacing/caps are threatened: execute the single most gate-decisive locomotion test first (\textbf{RUN} or \textbf{RUCK}, per phase contract) plus \textbf{CAL} battery; reslot the remaining locomotion test to Phase 05 Week 1 baseline window.
Never violate 96-hour maximal \textbf{RUN} vs maximal \textbf{RUCK} separation. \\

\end{longtblr}

\endgroup

\newpage

\subsection*{Phase 05 — Aerobic Base and Extended Stamina Development}
\phantomsection
\addcontentsline{toc}{subsubsection}{Phase 05 — Aerobic Base and Extended Stamina Development}

\begingroup
\scriptsize
\setlength{\tabcolsep}{2pt}
\renewcommand{\arraystretch}{1.12}

\begin{longtblr}{
  width=\linewidth,
  rowhead=1,
  colspec = {
    Q[l,wd=.7in]
    Q[l,wd=.80in]
    Q[l,wd=2.05in]
    X[l,1.25]
    X[l,3.50]
  },
  hlines,
  vlines,
  cells = {hsep=6pt, vsep=6pt, halign=l, valign=t},
  row{1} = {bg=StageBlue!15, fg=black, font=\bfseries, halign=l, valign=m},
  row{odd} = {bg=white},
  row{even} = {bg=LightGray},
}
\Hdr{Week-in-Phase} &
\Hdr{Week Range} &
\Hdr{Stress Type} &
\Hdr{Key Tests This Week (Major Tests Only)} &
\Hdr{Notes / Adjustments (Validity, Caps/Spacing, Environment, Reslot)} \\

Week 1 (baseline) &
Program W75 &
Deload+Test &
\textbf{CAL} battery;
\textbf{RUN} check-in: \textbf{C9} (\textbf{MET-2B-015}) (optional) &
Baseline window may include a conservative \textbf{RUN} check-in only if it does not compromise \textbf{CAL} battery validity.
Default is \textbf{CAL}-only if spacing is a concern. \\

Weeks 2–3 &
Program W76–W77 &
Build &
None (major tests not scheduled) &
Training-focused. \\

Week 4 &
Program W78 &
Deload+Test &
\textbf{RUN} \textbf{TT}: \textbf{C10} (\textbf{MET-2B-016})
\textbf{CAL} battery &
Two major events.
Enforce spacing. \\

Weeks 5–7 &
Program W79–W81 &
Intensify &
None (major tests not scheduled) &
Training-focused; preserve pool access planning and verification posture ahead of Week 8 if \textbf{SWIM} is scheduled. \\

Week 8 &
Program W82 &
Deload+Test &
\textbf{SWIM} \textbf{TT}: \textbf{C13} (\textbf{MET-2B-022} or \textbf{MET-2B-023}; unit locked)
\textbf{RUCK} \textbf{TT}: \textbf{C12} (\textbf{MET-2B-019}) &
Two major events.
\textbf{SWIM} requires \textbf{POOL}-2E-\#\#\# and verified unit/length; if pool is unavailable/unverifiable, reslot \textbf{SWIM} rather than substituting \textbf{AER}.
\textbf{RUCK} requires dry load verification and foot integrity.
Enforce spacing between \textbf{SWIM} and \textbf{RUCK} inside the week; if spacing cannot be preserved, prioritize the phase contract’s decisive event and reslot the other. \\

Weeks 9–11 &
Program W83–W85 &
Intensify &
None (major tests not scheduled) &
Training-focused; protect no-novelty posture approaching end-of-phase 5-mile \textbf{RUN} gate. \\

Week 12 &
Program W86 &
Deload+Test &
\textbf{RUN} \textbf{TT}: \textbf{C11} (\textbf{MET-2B-017}; 5-mile)
\textbf{CAL} battery (optional; cap-limited) &
End-of-phase gate anchored to high-cost \textbf{RUN}.
Default posture is \textbf{RUN} stand-alone (1 major event) or \textbf{RUN} + \textbf{CAL} (2 major events) if spacing allows and durability posture is stable.
Do not stack \textbf{RUCK} in the same week by default.
Any route/tool/footwear novelty is a comparability break and biases toward reslot for gate-grade use. \\

\end{longtblr}

\endgroup

\newpage

\subsection*{Phase 06 — Lactate Threshold and Power-Endurance Development}
\phantomsection
\addcontentsline{toc}{subsubsection}{Phase 06 — Lactate Threshold and Power-Endurance Development}

\begingroup
\scriptsize
\setlength{\tabcolsep}{2pt}
\renewcommand{\arraystretch}{1.12}

\begin{longtblr}{
  width=\linewidth,
  rowhead=1,
  colspec = {
    Q[l,wd=.7in]
    Q[l,wd=.80in]
    Q[l,wd=2.05in]
    X[l,1.25]
    X[l,3.50]
  },
  hlines,
  vlines,
  cells = {hsep=6pt, vsep=6pt, halign=l, valign=t},
  row{1} = {bg=StageBlue!15, fg=black, font=\bfseries, halign=l, valign=m},
  row{odd} = {bg=white},
  row{even} = {bg=LightGray},
}
\Hdr{Week-in-Phase} &
\Hdr{Week Range} &
\Hdr{Stress Type} &
\Hdr{Key Tests This Week (Major Tests Only)} &
\Hdr{Notes / Adjustments (Validity, Caps/Spacing, Environment, Reslot)} \\

Week 1 (baseline) &
Program W87 &
Deload+Test &
\textbf{CAL} battery (optional);
\textbf{RUN} check-in: \textbf{C8} (optional) &
Baseline is conservative.
Keep major event count minimal; do not insert multiple high-cost tests. \\

Weeks 2–4 &
Program W88–W90 &
Build &
None (major tests not scheduled) &
Training-focused. \\

Week 5 &
Program W91 &
Deload+Test &
\textbf{RUN} \textbf{TT}: \textbf{C8} (\textbf{MET-2B-014})
\textbf{CAL} battery &
Two major events.
Enforce spacing and preserve no-novelty posture. \\

Weeks 6–9 &
Program W92–W95 &
Intensify &
None (major tests not scheduled) &
Training-focused; avoid novelty in pool/machines approaching Week 10. \\

Week 10 &
Program W96 &
Deload+Test &
\textbf{RUCK} \textbf{TT}: \textbf{C12} (\textbf{MET-2B-019})
\textbf{SWIM} \textbf{TT}: \textbf{C13} (unit locked) &
Two major events.
Enforce spacing.
If pool access is compromised, reslot \textbf{SWIM}.
If foot integrity is compromised, reslot \textbf{RUCK}.
Do not substitute \textbf{SWIM} with \textbf{AER} for gate purposes. \\

Week 11 (taper intent) &
Program W97 &
Transition &
None (major tests not scheduled) &
Taper-aware stabilization week.
No new major tests.
Use only low-cost monitoring and readiness constraints to decide whether Week 12 can be executed gate-grade or must become recovery-first with reslot. \\

Week 12 (end gate) &
Program W98 &
Deload+Test &
\textbf{RUN} \textbf{TT}: single locked phase-contract variant only;
\textbf{CAL} battery &
End-of-phase gate week.
Keep to 1–2 major events.
If replication gaps exist, reslot rather than inserting a third major test that violates spacing. \\

\end{longtblr}

\endgroup

\newpage

\subsection*{Phase 07 — Advanced Hybrid Overload}
\phantomsection
\addcontentsline{toc}{subsubsection}{Phase 07 — Advanced Hybrid Overload}

\begingroup
\scriptsize
\setlength{\tabcolsep}{2pt}
\renewcommand{\arraystretch}{1.12}

\begin{longtblr}{
  width=\linewidth,
  rowhead=1,
  colspec = {
    Q[l,wd=.7in]
    Q[l,wd=.80in]
    Q[l,wd=2.05in]
    X[l,1.25]
    X[l,3.50]
  },
  hlines,
  vlines,
  cells = {hsep=6pt, vsep=6pt, halign=l, valign=t},
  row{1} = {bg=StageBlue!15, fg=black, font=\bfseries, halign=l, valign=m},
  row{odd} = {bg=white},
  row{even} = {bg=LightGray},
}
\Hdr{Week-in-Phase} &
\Hdr{Week Range} &
\Hdr{Stress Type} &
\Hdr{Key Tests This Week (Major Tests Only)} &
\Hdr{Notes / Adjustments (Validity, Caps/Spacing, Environment, Reslot)} \\

Week 1 (baseline) &
Program W99 &
Deload+Test &
\textbf{CAL} battery (optional) &
Conservative baseline. \\

Weeks 2–4 &
Program W100–W102 &
Build &
None &
Training-focused. \\

Week 5 &
Program W103 &
Deload+Test &
\textbf{STR} proxy session: \textbf{C17}, only if finalized and deployable;
\textbf{CAL} battery &
Two major events.
If spacing cannot be preserved inside the week, reslot \textbf{STR} to Week 10 or Week 12 (as phase needs permit). \\

Weeks 6–9 &
Program W104–W107 &
Intensify &
None &
Training-focused; preserve route/pool stability approaching Week 10. \\

Week 10 &
Program W108 &
Deload+Test &
\textbf{RUN} \textbf{TT}: \textbf{C10} (\textbf{MET-2B-016})
\textbf{SWIM} \textbf{TT}: \textbf{C13} &
Two major events.
Enforce spacing. \\

Week 11 (taper intent) &
Program W109 &
Transition &
None &
Stabilization week.
No new major tests. \\

Week 12 (end gate) &
Program W110 &
Deload+Test &
\textbf{RUCK} \textbf{TT}: \textbf{C12} (\textbf{MET-2B-019})
\textbf{CAL} battery &
End-of-phase gate week.
Enforce spacing.
If durability flags are active, reslot \textbf{RUCK} rather than forcing a compromised attempt. \\

\end{longtblr}

\endgroup

\newpage

\subsection*{Phase 08 — Structural Durability and Load Tolerance}
\phantomsection
\addcontentsline{toc}{subsubsection}{Phase 08 — Structural Durability and Load Tolerance}

\begingroup
\scriptsize
\setlength{\tabcolsep}{2pt}
\renewcommand{\arraystretch}{1.12}

\begin{longtblr}{
  width=\linewidth,
  rowhead=1,
  colspec = {
    Q[l,wd=.7in]
    Q[l,wd=.80in]
    Q[l,wd=2.05in]
    X[l,1.25]
    X[l,3.50]
  },
  hlines,
  vlines,
  cells = {hsep=6pt, vsep=6pt, halign=l, valign=t},
  row{1} = {bg=StageBlue!15, fg=black, font=\bfseries, halign=l, valign=m},
  row{odd} = {bg=white},
  row{even} = {bg=LightGray},
}
\Hdr{Week-in-Phase} &
\Hdr{Week Range} &
\Hdr{Stress Type} &
\Hdr{Key Tests This Week (Major Tests Only)} &
\Hdr{Notes / Adjustments (Validity, Caps/Spacing, Environment, Reslot)} \\

Week 1 (baseline) &
Program W111 &
Deload+Test &
\textbf{CAL} battery (optional)
\textbf{RUN} check-in (optional) &
Conservative baseline. \\

Weeks 2–4 &
Program W112–W114 &
Build &
None &
Training-focused. \\

Week 5 &
Program W115 &
Deload+Test &
\textbf{RUN} \textbf{TT}: \textbf{C10} (\textbf{MET-2B-016})
\textbf{CAL} battery &
Two major events.
Enforce spacing. \\

Weeks 6–9 &
Program W116–W119 &
Intensify &
None &
Training-focused; protect foot integrity ahead of long ruck milestone. \\

Week 10 &
Program W120 &
Deload+Test &
Long \textbf{RUCK} milestone: \textbf{C12} (\textbf{MET-2B-021}; 12 mi / 55 lb dry) (stand-alone by intent) &
Stand-alone major event by default (1 major event week).
Do not pair with maximal \textbf{RUN} in this week.
If the long ruck construct is invalid due to execution or logging defects, reslot rather than substituting.
Dry load verification required; water logged separately; environment hazards trigger reslot. \\

Week 11 (taper intent) &
Program W121 &
Transition &
None &
Stabilization week; no new major tests. \\

Week 12 (end gate) &
Program W122 &
Deload+Test &
\textbf{SWIM} \textbf{TT}: \textbf{C13}
\textbf{RUCK} \textbf{TT}: \textbf{C12} (\textbf{MET-2B-019}) (cap/spacing permitting) &
Keep to 1–2 major events.
If Week 10 executed long ruck, treat Week 12 as lower density and prioritize the most gate-decisive remaining domain.
If spacing/caps are threatened, reslot the lesser-priority event. \\

\end{longtblr}

\endgroup

\newpage

\subsection*{Phase 09 — High-Volume Calisthenics and Stamina}
\phantomsection
\addcontentsline{toc}{subsubsection}{Phase 09 — High-Volume Calisthenics and Stamina}

\begingroup
\scriptsize
\setlength{\tabcolsep}{2pt}
\renewcommand{\arraystretch}{1.12}

\begin{longtblr}{
  width=\linewidth,
  rowhead=1,
  colspec = {
    Q[l,wd=.7in]
    Q[l,wd=.80in]
    Q[l,wd=2.05in]
    X[l,1.25]
    X[l,3.50]
  },
  hlines,
  vlines,
  cells = {hsep=6pt, vsep=6pt, halign=l, valign=t},
  row{1} = {bg=StageBlue!15, fg=black, font=\bfseries, halign=l, valign=m},
  row{odd} = {bg=white},
  row{even} = {bg=LightGray},
}
\Hdr{Week-in-Phase} &
\Hdr{Week Range} &
\Hdr{Stress Type} &
\Hdr{Key Tests This Week (Major Tests Only)} &
\Hdr{Notes / Adjustments (Validity, Caps/Spacing, Environment, Reslot)} \\

Week 1 (baseline) &
Program W123 &
Deload+Test &
\textbf{CAL} battery (central)
\textbf{RUN} check-in optional &
Baseline emphasizes \textbf{CAL} execution quality and artifact compliance. \\

Weeks 2–4 &
Program W124–W126 &
Build &
None &
Training-focused; do not introduce gate-grade tests outside protected windows. \\

Week 5 &
Program W127 &
Deload+Test &
\textbf{CAL} battery
\textbf{RUN} \textbf{TT}: \textbf{C9} (optional if replication needed) &
Keep density controlled: \textbf{CAL} battery is primary; add \textbf{RUN} only when replication gaps exist and spacing can be preserved. \\

Weeks 6–9 &
Program W128–W131 &
Intensify &
None &
Training-focused; monitor durability flags closely due to high calisthenics volume. \\

Week 10 &
Program W132 &
Deload+Test &
\textbf{SWIM} \textbf{TT}: \textbf{C13}
\textbf{RUCK} \textbf{TT}: \textbf{C12} (\textbf{MET-2B-019}) &
Two major events.
Enforce spacing.
Reslot if pool access fails or foot integrity fails. \\

Week 11 (taper intent) &
Program W133 &
Transition &
None &
Stabilization week; no new major tests. \\

Week 12 (end gate) &
Program W134 &
Deload+Test &
\textbf{RUN} \textbf{TT}: \textbf{C11} (\textbf{MET-2B-017})
\textbf{CAL} battery &
High-cost \textbf{RUN}.
Keep to 1–2 major events; default is \textbf{RUN} + \textbf{CAL} only.
Do not add ruck in this week by default. \\

\end{longtblr}

\endgroup

\newpage

\subsection*{Phase 10 — Advanced Tactical Readiness}
\phantomsection
\addcontentsline{toc}{subsubsection}{Phase 10 — Advanced Tactical Readiness}

\begingroup
\scriptsize
\setlength{\tabcolsep}{2pt}
\renewcommand{\arraystretch}{1.12}

\begin{longtblr}{
  width=\linewidth,
  rowhead=1,
  colspec = {
    Q[l,wd=.7in]
    Q[l,wd=.80in]
    Q[l,wd=2.05in]
    X[l,1.25]
    X[l,3.50]
  },
  hlines,
  vlines,
  cells = {hsep=6pt, vsep=6pt, halign=l, valign=t},
  row{1} = {bg=StageBlue!15, fg=black, font=\bfseries, halign=l, valign=m},
  row{odd} = {bg=white},
  row{even} = {bg=LightGray},
}
\Hdr{Week-in-Phase} &
\Hdr{Week Range} &
\Hdr{Stress Type} &
\Hdr{Key Tests This Week (Major Tests Only)} &
\Hdr{Notes / Adjustments (Validity, Caps/Spacing, Environment, Reslot)} \\

Week 1 (baseline) &
Program W135 &
Deload+Test &
\textbf{CAL} battery
\textbf{RUN} check-in: \textbf{C8} (optional) &
Baseline remains conservative. \\

Weeks 2–4 &
Program W136–W138 &
Build &
None &
Training-focused; protect pool governance prerequisites if underwater is ever expected later. \\

Week 5 &
Program W139 &
Deload+Test &
\textbf{SWIM} \textbf{TT}: \textbf{C13}
\textbf{CAL} battery &
Two major events.
Enforce spacing. \\

Weeks 6–9 &
Program W140–W143 &
Intensify &
None &
Training-focused; strict no-novelty posture approaching Week 10–12 windows. \\

Week 10 &
Program W144 &
Deload+Test &
\textbf{RUN} \textbf{TT}: \textbf{C10} (\textbf{MET-2B-016})
Underwater governed test: \textbf{C14} (\textbf{MET-2B-025}; \textbf{MET-2B-026}) conditional &
If underwater prerequisites are not fully satisfied and logged, underwater is not scheduled.
If underwater is scheduled, keep the week at 1–2 major events and do not schedule long ruck in the same protected week.
Enforce spacing and treat any governance ambiguity as a reslot trigger. \\

Week 11 (taper intent) &
Program W145 &
Transition &
None &
Stabilization week; no new major tests. \\

Week 12 (end gate) &
Program W146 &
Deload+Test &
\textbf{RUN} \textbf{TT}: \textbf{C11} (\textbf{MET-2B-017})
\textbf{RUCK} \textbf{TT}: \textbf{C12} (\textbf{MET-2B-019}) (staged; cap/spacing permitting) &
Two major events.
Enforce 96-hour separation if maximal variants are involved.
If spacing cannot be preserved inside the week, prioritize the dominant gate objective and reslot the other test to Phase 11 Week 1 baseline or the next protected window.
Underwater is not added to this week by default. \\

\end{longtblr}

\endgroup

\newpage

\subsection*{Phase 11 — Pre-Selection Conditioning and Recovery}
\phantomsection
\addcontentsline{toc}{subsubsection}{Phase 11 — Pre-Selection Conditioning and Recovery}

\begingroup
\scriptsize
\setlength{\tabcolsep}{2pt}
\renewcommand{\arraystretch}{1.12}

\begin{longtblr}{
  width=\linewidth,
  rowhead=1,
  colspec = {
    Q[l,wd=.7in]
    Q[l,wd=.80in]
    Q[l,wd=2.05in]
    X[l,1.25]
    X[l,3.50]
  },
  hlines,
  vlines,
  cells = {hsep=6pt, vsep=6pt, halign=l, valign=t},
  row{1} = {bg=StageBlue!15, fg=black, font=\bfseries, halign=l, valign=m},
  row{odd} = {bg=white},
  row{even} = {bg=LightGray},
}
\Hdr{Week-in-Phase} &
\Hdr{Week Range} &
\Hdr{Stress Type} &
\Hdr{Key Tests This Week (Major Tests Only)} &
\Hdr{Notes / Adjustments (Validity, Caps/Spacing, Environment, Reslot)} \\

Week 1 (baseline) &
Program W147 &
Deload+Test &
\textbf{CAL} battery (optional) &
Conservative baseline; recovery-sensitive phase. \\

Weeks 2–4 &
Program W148–W150 &
Build &
None &
Training-focused. \\

Week 5 &
Program W151 &
Deload+Test &
\textbf{RUN} \textbf{TT}: \textbf{C10} (\textbf{MET-2B-016})
\textbf{CAL} battery &
Two major events.
Enforce spacing. \\

Weeks 6–9 &
Program W152–W155 &
Intensify &
None &
Training-focused; maintain configuration stability approaching Week 10–12. \\

Week 10 &
Program W156 &
Deload+Test &
Long \textbf{RUCK} milestone: \textbf{C12} (\textbf{MET-2B-021}) (stand-alone)
or
\textbf{SWIM} \textbf{TT} \textbf{C13} + \textbf{RUCK} \textbf{TT} \textbf{MET-2B-019} &
If long ruck is scheduled, treat it as stand-alone major event by default.
If long ruck is not scheduled, the alternative is two major events (\textbf{SWIM} + 6-mile ruck) under spacing rules.
Any foot/gait flags force reslot. \\

Week 11 (taper intent) &
Program W157 &
Transition &
None &
Stabilization week; no new major tests. \\

Week 12 (end gate) &
Program W158 &
Deload+Test &
\textbf{RUN} \textbf{TT}: \textbf{C11} (\textbf{MET-2B-017})
\textbf{CAL} battery &
High-cost \textbf{RUN}.
Keep to 1–2 major events.
Underwater is optional only under prerequisites and is not scheduled by default. \\

\end{longtblr}

\endgroup

\newpage

\subsection*{Phase 12 — Final Pre-Selection Conditioning}
\phantomsection
\addcontentsline{toc}{subsubsection}{Phase 12 — Final Pre-Selection Conditioning}

\begingroup
\scriptsize
\setlength{\tabcolsep}{2pt}
\renewcommand{\arraystretch}{1.12}

\begin{longtblr}{
  width=\linewidth,
  rowhead=1,
  colspec = {
    Q[l,wd=.7in]
    Q[l,wd=.80in]
    Q[l,wd=2.05in]
    X[l,1.25]
    X[l,3.50]
  },
  hlines,
  vlines,
  cells = {hsep=6pt, vsep=6pt, halign=l, valign=t},
  row{1} = {bg=StageBlue!15, fg=black, font=\bfseries, halign=l, valign=m},
  row{odd} = {bg=white},
  row{even} = {bg=LightGray},
}
\Hdr{Week-in-Phase} &
\Hdr{Week Range} &
\Hdr{Stress Type} &
\Hdr{Key Tests This Week (Major Tests Only)} &
\Hdr{Notes / Adjustments (Validity, Caps/Spacing, Environment, Reslot)} \\

Week 1 (baseline) &
Program W159 &
Deload+Test &
\textbf{CAL} battery (optional);
\textbf{RUN} check-in: optional &
Conservative baseline; begin configuration-freeze posture for routes/tools used in final verification windows. \\

Weeks 2–4 &
Program W160–W162 &
Build &
None &
Training-focused; no gate-grade time trials. \\

Week 5 &
Program W163 &
Deload+Test &
\textbf{RUN} \textbf{TT}: \textbf{C9} (\textbf{MET-2B-015})
\textbf{CAL} battery &
Two major events.
Enforce spacing. \\

Weeks 6–8 &
Program W164–W166 &
Intensify &
None &
Training-focused; strict no-novelty posture approaching verification block. \\

Weeks 9–10 (verification block) &
Program W167–W168 &
Deload+Test (two-week verification block) &
Targeted retests only: \textbf{RUN} (\textbf{C8}–\textbf{C11}), \textbf{RUCK} (\textbf{C12}), \textbf{SWIM} (\textbf{C13}), \textbf{STR} (\textbf{C17}), \textbf{AER} (\textbf{C18}) as applicable; staged to satisfy replication without stacking &
Treat as two consecutive protected windows with spacing enforced across both weeks.
Do not attempt to execute all domains; prioritize weakest domains needing replication (2 of last 3 \textbf{VALID}).
If a comparability break occurs (new route/machine/pool), treat outcomes conservatively and extend verification rather than compressing.
Underwater (\textbf{C14}) is optional only under prerequisites and is not scheduled by default. \\

Week 11 (taper intent) &
Program W169 &
Transition &
None &
Stabilization week; no new major tests. \\

Week 12 (final gate) &
Program W170 &
Deload+Test &
Final gate selection: \textbf{CAL} battery + \textbf{RUN} \textbf{TT} \textbf{C11}
or
long \textbf{RUCK} milestone \textbf{C12} (\textbf{MET-2B-021}) &
Choose dominant gate need; do not stack all maximal domains.
Maintain cap discipline (1–2 major events).
If the dominant event is long ruck, keep the week otherwise minimal.
If replication gaps remain after W9–10, reslot remaining gaps rather than adding extra major events that violate spacing. \\

\end{longtblr}

\endgroup

\newpage

\subsection*{Phase 13 — Full-Spectrum Readiness Consolidation}
\phantomsection
\addcontentsline{toc}{subsubsection}{Phase 13 — Full-Spectrum Readiness Consolidation}

\begingroup
\scriptsize
\setlength{\tabcolsep}{2pt}
\renewcommand{\arraystretch}{1.12}

\begin{longtblr}{
  width=\linewidth,
  rowhead=1,
  colspec = {
    Q[l,wd=.7in]
    Q[l,wd=.80in]
    Q[l,wd=2.05in]
    X[l,1.25]
    X[l,3.50]
  },
  hlines,
  vlines,
  cells = {hsep=6pt, vsep=6pt, halign=l, valign=t},
  row{1} = {bg=StageBlue!15, fg=black, font=\bfseries, halign=l, valign=m},
  row{odd} = {bg=white},
  row{even} = {bg=LightGray},
}
\Hdr{Week-in-Phase} &
\Hdr{Week Range} &
\Hdr{Stress Type} &
\Hdr{Key Tests This Week (Major Tests Only)} &
\Hdr{Notes / Adjustments (Validity, Caps/Spacing, Environment, Reslot)} \\

Week 1 (baseline) &
Program W171 &
Deload+Test &
\textbf{CAL} battery (optional)
\textbf{BODYCOMP} trend capture continues &
Conservative baseline; confirm configuration freeze posture for the final verification cycle (routes, pool unit/length, treadmill IDs, footwear). \\

Weeks 2–4 &
Program W172–W174 &
Build &
None &
Training-focused; no gate-grade time trials. \\

Week 5 &
Program W175 &
Deload+Test &
\textbf{SWIM} \textbf{TT}: \textbf{C13} (unit locked)
\textbf{CAL} battery &
Two major events.
Enforce spacing.
Pool verification is mandatory; if pool access is compromised, reslot \textbf{SWIM}. \\

Weeks 6–8 &
Program W176–W178 &
Intensify &
None &
Training-focused; preserve stability and avoid novelty. \\

Weeks 9–10 (targeted retest block) &
Program W179–W180 &
Deload+Test (two-week retest block) &
Targeted weakest-domain retests only: \textbf{RUN} \textbf{C8}–\textbf{C11}; \textbf{RUCK} \textbf{C12}; \textbf{SWIM} \textbf{C13}; \textbf{STR} \textbf{C17}; \textbf{AER} \textbf{C18} as applicable; staged for replication completion &
Treat as two consecutive protected windows with spacing enforced across both weeks.
Objective is replication closure (2 of last 3 \textbf{VALID}) under stable configurations.
Do not schedule “simulation” constructs outside Stage 2 protocols.
Underwater remains optional only under prerequisites; if unsupported it is not scheduled and cannot be claimed via proxy. \\

Week 11 (taper intent) &
Program W181 &
Transition &
None &
Stabilization week; no new major tests. \\

Week 12 (final verification) &
Program W182 &
Deload+Test &
Final verification selection: prioritized \textbf{RUN} \textbf{TT} (\textbf{C8}–\textbf{C11}) + \textbf{CAL} battery
or
long \textbf{RUCK} milestone \textbf{C12} (\textbf{MET-2B-021}) as dominant event &
Keep the week minimal and decisive.
Do not stack all maximal domains.
If the dominant event is long ruck, do not add maximal \textbf{RUN} in the same week; preserve spacing and validity.
If final replication gaps remain, default posture is \textbf{HOLD} (Replication) for the affected metric rather than forcing compromised evidence. \\

\end{longtblr}

\endgroup

\end{landscape}

\subsection*{Calendar Integrity Check}
\phantomsection
\addcontentsline{toc}{subsection}{Calendar Integrity Check}

\begin{itemize}
\item Every phase includes at least one mid-phase protected Deload+Test window and an end-of-phase Deload+Test exit gate window by design.
\item Week 1 (baseline) exists in every phase and is treated as an infrastructure and standards touchpoint rather than a maximal gate posture by default.
\item Major tests are scheduled only in Deload+Test weeks (protected windows) or in explicitly defined two-week verification blocks (Phase 12–13 Week 9–10). Build/Intensify/Transition weeks do not schedule gate-grade major tests by default.
\item When multiple domains appear in the same Deload+Test week, the Notes/Adjustments field encodes staged battery intent and prioritization consistent with cap and spacing doctrine.
\item \textbf{CAL} gate-grade validity requires compliant artifact evidence per Stage 2.C; missing/non-compliant artifacts force \textbf{INVALID} for gate use and trigger reslot posture rather than stacking immediate reattempts.
\item \textbf{RUN} and \textbf{RUCK} scheduling respects the separation intent: maximal \textbf{RUN} and maximal \textbf{RUCK} are not co-located within the same 96-hour window; long ruck milestones are treated as stand-alone major events by default.
\item \textbf{SWIM} unit handling is explicit and unit locked: 500-yard and 500-meter tests are separate evidence streams; no conversion is implied for gate claims. \textbf{POOL}-2E-\#\#\# verification is mandatory.
\item Underwater governance posture is explicit: underwater appears only as conditional, supervised-only \textbf{C14} constructs in Phase 10–13, and is explicitly “not scheduled” when prerequisites cannot be guaranteed.
\item Environment volatility and tool drift are treated as predictable constraints: unsafe conditions or comparability breaks trigger reslot posture or indoor substitution only when Stage 2 validity/logging preserves comparability.
\item Replication posture is supported without overtesting: two-week verification blocks exist in Phases 12–13 to close replication gaps across windows while preserving cap and spacing discipline.
\end{itemize}

\end{document}